\scanpage{38}
% archive sheet number: 37
\deleted{Soit} $\mathcal X$ ens. ordonné fini (on ne suppose pas qu'il
ait un plus petit ou plus grand élément). On
note \inserted{$|\mathcal X|=$} réal. $(\mathcal X)$ sa réalisation topologique (avec sa
décomposition \unsure{cônique|conique} canonique, permettant
d'identifier $\mathcal X$ à un ens ordonné de parties
fermées de $\mathcal X$ \inserted{recouvrant $\mathcal X$}). Pour tt $x\in\mathcal X$, soit
$\mathcal X_{<x}=$ \deleted{\unsure{l'ens des} $y\in$} $\{y\mid y\in\mathcal X,\ y<x\}$. Faisons l'hypothèse :

\[
\forall x\in\mathcal X,\qquad |\mathcal X_{<x}|\ \text{est une sphère topologique.}\tag*{\llap{\makebox[\displaywidth][l]{(R)}}}
\]

Soit $\delta(x)=\dim|\mathcal X_{<x}|+1=\dim|\mathcal X_{\leq x}|$. On trouve

\uline{Proposition} Si $x<y$, on a $\delta(x)<\delta(y)$ (i.e. $\delta:\mathcal X\to\mathbb N$
est \unsure{strictmt|strictement} croissante). Pour qu'on ait \deleted{$\delta(y)$}
\[
\delta(y)=\delta(x)+1
\]
il f. et s. que $y$ soit un successeur de $x$.

Première assertion claire, \deleted{Pour} compte tenu
que $\dim\underbrace{C(X)}_{\text{cône sur }X}=\dim X+1$, et dim est
une fonction croissante\ldots{} L'assertion ``il faut''
en résulte, prouvons ``il suffit''. Supposons que
$y$ soit un successeur de $x$. Considérons
$S=|\mathcal X_{<y}|$ sphère de dim $d\unsure{=|:}\delta(y)-1\geq\delta(x)$,
contenant $|\mathcal X_{\leq x}|$, qui est une facette
% slanted
\margin{vérifier en général sur une variété}
maximale, donc de dim $=\deleted{\unsure{\delta}}d$, donc $\delta(x)=\delta(y)-1$,
cqfd. [Pour qu'on ait $\delta(y)=\delta(x)+1$, il suffit que
$y$ soit succ. de $x$ et $|\mathcal X_{<y}|$ une \uline{variété} topologique.
